\begin{table}[htbp]
\centering
\caption{Bajari-Ye Placebo: Random Group Assignment}
\label{tab:bajari_ye_placebo}
\begin{adjustbox}{max width=\textwidth}
\begin{threeparttable}
\small
\begin{tabular}{lcc}
\toprule
 & Mean Pairwise Product & p-value \\
\midrule
FL pairs (observed) & 5.1597 & 0.000000 \\
Non-FL pairs (observed) & 2.2082 & 0.000000 \\
Fake random groups (placebo) & 5.8029 & 0.0000 \\
\midrule
Difference (FL $-$ non-FL) & 2.9515 & \\
Bootstrap SE & 0.6839 & \\
Bootstrap 95\% CI & [2.5675, 5.2094] & \\
\bottomrule
\end{tabular}
\begin{tablenotes}
\small
\item \textit{Notes:} Pairwise products of first-stage bid residuals % V7-FIX3: placebo note rewritten
within tenders. The fake-groups placebo demonstrates that tender-level
shocks induce positive residual correlations among all bidders in the
same tender. The fake-groups product (5.80) exceeds zero for the same
reason the non-FL product (2.21) exceeds zero: common cost signals
within the tender. The economically relevant statistic is the
FL\,--\,non-FL pairwise difference (2.95, bootstrap 95\% CI:
$[2.57,\; 5.21]$), which isolates excess within-tender correlation
of FL bids above and beyond common shocks shared by all bidders.
Bootstrap: 1,000 iterations resampling tenders with replacement.
\end{tablenotes}
\end{threeparttable}
\end{adjustbox}
\end{table}
